\honest{The Sacred Mundane}{Eliezer Yudkowsky}{2009}

Before a video dialogue with Adam Frank, I read about the first half of his book \textsc{The Constant Fire}. It is about the experience of the sacred, which I know from watching a shuttle launch, or the stars. Frank holds that science and religion share this experience. I hold that it is a human quality that religion corrupts.

Frank quotes William James: religion ``shall mean for us the feelings, acts, and experiences of individual men in their solitude.'' The book goes on to say that the sacred is private and individual. \nb{James offered the definition ``arbitrarily,'' as a working limit for lectures that would leave churches and theology aside. From Frank's own text the post quotes only ``unique experience of the sacred.''}

This baffled me. Must I feel nothing sacred because many others watch the same launch? If ``private'' only means hard to put into words, why stress it of the sacred rather than of sneezing? Then I saw it: a private experience cannot be criticized. And the price of that shield is solitude, which James admired ``as if loneliness were a good thing.'' \nb{James's solitude is the individual as against the church. He put personal religion first because the founders of churches drew their power from ``direct personal communion with the divine.'' Neither James nor Frank is shown using privacy to fend off a critic; the line of defense the post quotes is its own.}

Religion twists the sacred, I say, through five habits that grew up to shield its lie. Mysteriousness: ``It is a sacred mystery!'' answers hard questions. Faith: when the evidence did not come, religion made ``the retreat to commitment.'' \nb{This passage repeats, nearly word for word, one in Yudkowsky's ``Religion's Claim to be Non-Disprovable'' (2007).} The retreat to pure experience, which strips the feeling of its object. Separation, which makes human works like the shuttle seem unfit to be sacred, as Keats demoted the rainbow to the ``dull catalogue of mundane things.'' \nb{Keats wrote ``common things''; ``mundane'' is this post's word.} And privacy. For each I name a cost.

So we had best not salvage religion, not even as ``spirituality''. Without the churches and scriptures, what is left is the five habits. \nb{The post examines no spiritual practice, and the one book it discusses is faulted only for its stress on privacy.} The habits evolved to defend a lie; spirituality is the poisoned cup with the pellet taken out. The wise, I think, start over from scratch. \nb{In ``The Genetic Fallacy'' (2008) Yudkowsky had endorsed doubting everything that grew from a discredited source, and added: ``You are not licensed to reject them outright as conclusions.'' The costs the post names for each habit are reasons that do not depend on origin.}

The capacity to admit you were entirely wrong is why religious experience will never be like scientific experience. I end with the stars: believable, real, knowable, shared, made of the same fabric as everything else, ``the sacred mundane.''
